\documentclass[11pt]{article}

\usepackage{amsmath,amsthm,amssymb,mathtools,bbm,enumitem,geometry}
\geometry{margin=1in}

\newcommand{\R}{\mathbb{R}}
\newcommand{\1}{\mathbf{1}}

\newtheorem{theorem}{Theorem}
\newtheorem{lemma}{Lemma}
\newtheorem{corollary}{Corollary}
\newtheorem{remark}{Remark}

\title{Asymptotic Upper Bound of a Transfer Learning Estimator}


\begin{document}
\section{Results from the zero-shot prompt.}
\paragraph{Notation and interpretation of the risk.}
Let
\[
\mathcal R:=\mathcal R(t)\subseteq [p],\qquad q:=|\mathcal R|,\qquad \mathcal U:=\mathcal R^c,
\]
and write the covariate as
\[
x=(z^\top,u^\top)^\top,\qquad z:=x_{\mathcal R}\in\mathbb R^q,\quad u:=x_{\mathcal U}\in\mathbb R^{p-q}.
\]
Partition the covariance matrix as
\[
\Sigma=
\begin{pmatrix}
\Sigma_{RR}&\Sigma_{RU}\\
\Sigma_{UR}&\Sigma_{UU}
\end{pmatrix},
\qquad \Sigma_{RR}\in\mathbb R^{q\times q}.
\]
For the target site, write
\[
\beta^{(t)}=
\begin{pmatrix}
\beta_R^{(t)}\\
\beta_U^{(t)}
\end{pmatrix},
\qquad
\beta^{(s)}=
\begin{pmatrix}
\beta_R^{(s)}\\
\beta_U^{(s)}
\end{pmatrix}.
\]

The estimator in the statement should use the available source estimator \(\hat\beta^{(s)}\); accordingly we interpret
\[
\hat\omega_s^{(t)}=\Phi^{(t)}\hat\beta^{(s)},
\qquad
\Phi^{(t)}=\bigl[I_q,\ \Sigma_{RR}^{-1}\Sigma_{RU}\bigr].
\]
This is the natural transfer center induced by the source estimate.

Since \(\hat\omega_{\mathrm{TL}}^{(t)}\) depends on the target responses \(Y^{(t)}\) and on \(\hat\beta^{(s)}\), the quantity
\[
r^{\mathrm{pred}}\!\bigl(\hat\omega_{\mathrm{TL}}^{(t)}\bigr)
=
\mathbb E\!\left[
\bigl(y^{(t)}-z^\top\hat\omega_{\mathrm{TL}}^{(t)}\bigr)^2
\,\middle|\,
X_R^{(t)}
\right]
\]
is still random through the training responses and the source estimator. The standard object in random-design asymptotics is therefore its conditional mean
\[
\bar r^{\mathrm{pred}}(\lambda_\omega\mid X_R^{(t)})
:=
\mathbb E\!\left[
r^{\mathrm{pred}}\!\bigl(\hat\omega_{\mathrm{TL}}^{(t)}\bigr)
\,\middle|\,
X_R^{(t)}
\right],
\]
where the outer expectation is over the remaining randomness
(\(Y^{(t)}\), \(\beta^{(t)}\), \(\beta^{(s)}\), \(\hat\beta^{(s)}\)).
We shall derive a finite-sample exact decomposition for \(\bar r^{\mathrm{pred}}\), then a high-dimensional asymptotic upper bound, and then the oracle choice of \(\lambda_\omega\).

\bigskip

\paragraph{Assumptions.}
We use the following regularity conditions.

\begin{enumerate}
\item[(A1)] \textbf{High-dimensional regime.}
As \(n_t,p\to\infty\),
\[
\frac{p}{n_t}\to\gamma\in(0,\infty),
\qquad
\frac{q}{p}\to\kappa\in(0,1],
\qquad
\frac{q}{n_t}\to\eta:=\kappa\gamma.
\]

\item[(A2)] \textbf{Covariance regularity.}
The matrix \(\Sigma\) is positive definite, and there exist constants
\(0<c_\Sigma<C_\Sigma<\infty\) such that for all \(p\),
\[
c_\Sigma\le \lambda_{\min}(\Sigma)\le \lambda_{\max}(\Sigma)\le C_\Sigma.
\]
Hence also
\[
c_\Sigma\le \lambda_{\min}(\Sigma_{RR})\le \lambda_{\max}(\Sigma_{RR})\le C_\Sigma.
\]

\item[(A3)] \textbf{Transfer map regularity.}
There exists a finite constant \(C_\Phi\) such that
\[
\|\Phi^{(t)}\|_{\mathrm{op}}^2
=
\lambda_{\max}\!\bigl(\Phi^{(t)}(\Phi^{(t)})^\top\bigr)
\le C_\Phi
\qquad\text{for all }p.
\]

\item[(A4)] \textbf{Random coefficients.}
\[
\mathbb E[\beta^{(m)}]=0,\qquad
\mathrm{Var}(\beta^{(m)})=\frac{\alpha_m^2}{p}I_p,\qquad m\in\{s,t\},
\]
and
\[
\mathrm{Cov}\!\bigl(\beta^{(t)},\beta^{(s)}\bigr)
=
\frac{\rho\alpha_t\alpha_s}{p}I_p,
\qquad -1\le \rho\le 1.
\]

\item[(A5)] \textbf{Source estimation error.}
Let
\[
e_s:=\hat\beta^{(s)}-\beta^{(s)}.
\]
Assume \(e_s\) is independent of the target sample \((X^{(t)},Y^{(t)})\) and of \(\beta^{(t)}\), satisfies \(\mathbb E[e_s]=0\), and
\[
\lambda_{\max}\!\left(\mathbb E[e_s e_s^\top]\right)\le \frac{C^U}{p}
\]
for a constant \(C^U>0\).

\item[(A6)] \textbf{Noise.}
The target noise variables \(\varepsilon_i\) are i.i.d., mean zero, variance \(\sigma_\varepsilon^2\), and independent of \(x_i\), \(\beta^{(t)}\), \(\beta^{(s)}\), and \(e_s\).

\item[(A7)] \textbf{Spectral convergence of the observed design.}
The empirical spectral distribution of
\[
S_n:=\frac1{n_t}(X_R^{(t)})^\top X_R^{(t)}
\]
converges almost surely to a deterministic compactly supported probability law \(F_\eta\).
This is automatic, for example, if \(z_i\sim N(0,\Sigma_{RR})\) and the empirical spectral distribution of \(\Sigma_{RR}\) converges, in which case \(F_\eta\) is the corresponding deformed Marchenko--Pastur limit.
\end{enumerate}

\bigskip

\paragraph{Step 1: the target optimal linear predictor from \(z=x_{\mathcal R}\).}

\begin{lemma}[Optimal observed-feature coefficient]
Under the Gaussian model,
\[
\omega_*^{(t)}:=\arg\min_{\omega\in\mathbb R^q}
\mathbb E\!\left[(y^{(t)}-z^\top\omega)^2\,\middle|\,\beta^{(t)}\right]
=
\Phi^{(t)}\beta^{(t)}
=
\beta_R^{(t)}+\Sigma_{RR}^{-1}\Sigma_{RU}\beta_U^{(t)}.
\]
Moreover, writing
\[
\Sigma_{U\mid R}:=\Sigma_{UU}-\Sigma_{UR}\Sigma_{RR}^{-1}\Sigma_{RU},
\]
we have the conditional representation
\[
y^{(t)}=z^\top\omega_*^{(t)}+\xi,
\]
where \(\xi\) is independent of \(z\) given \(\beta^{(t)}\), with
\[
\mathbb E[\xi\mid \beta^{(t)}]=0,
\qquad
\mathrm{Var}(\xi\mid \beta^{(t)})
=
\tau_t^2(\beta^{(t)})
:=
\sigma_\varepsilon^2+(\beta_U^{(t)})^\top\Sigma_{U\mid R}\beta_U^{(t)}.
\]
\end{lemma}

\begin{proof}
Since \((z,u)\) is jointly Gaussian with mean \(0\),
\[
\mathbb E[u\mid z]=\Sigma_{UR}\Sigma_{RR}^{-1}z.
\]
Therefore
\[
\mathbb E[y^{(t)}\mid z,\beta^{(t)}]
=
z^\top\beta_R^{(t)}+\mathbb E[u\mid z]^\top\beta_U^{(t)}
=
z^\top\Bigl(\beta_R^{(t)}+\Sigma_{RR}^{-1}\Sigma_{RU}\beta_U^{(t)}\Bigr).
\]
The coefficient of the conditional mean is therefore \(\omega_*^{(t)}=\Phi^{(t)}\beta^{(t)}\). The decomposition
\[
\xi:=y^{(t)}-z^\top\omega_*^{(t)}
=
u^\top\beta_U^{(t)}-\mathbb E[u^\top\beta_U^{(t)}\mid z]+\varepsilon
\]
has conditional mean zero. Since, conditional on \(\beta^{(t)}\), the Gaussian residual \(u-\mathbb E[u\mid z]\) is independent of \(z\) and has covariance \(\Sigma_{U\mid R}\),
\[
\mathrm{Var}(\xi\mid \beta^{(t)})
=
(\beta_U^{(t)})^\top\Sigma_{U\mid R}\beta_U^{(t)}+\sigma_\varepsilon^2.
\]
This proves the claim.
\end{proof}

\bigskip

\paragraph{Step 2: closed form of the transfer-ridge estimator.}

Set
\[
X:=X_R^{(t)}\in\mathbb R^{n_t\times q},\qquad
Y:=Y^{(t)}\in\mathbb R^{n_t},\qquad
S_n:=\frac1{n_t}X^\top X.
\]
The estimator is
\[
\hat\omega_{\mathrm{TL}}
=
\arg\min_{\omega\in\mathbb R^q}
\left[
\frac1{n_t}\|Y-X\omega\|_2^2+\lambda_\omega\|\omega-\hat\omega_s\|_2^2
\right].
\]
Its first-order condition gives
\[
(S_n+\lambda_\omega I_q)\hat\omega_{\mathrm{TL}}
=
\frac1{n_t}X^\top Y+\lambda_\omega \hat\omega_s,
\]
hence
\[
\hat\omega_{\mathrm{TL}}
=
(S_n+\lambda_\omega I_q)^{-1}
\left(\frac1{n_t}X^\top Y+\lambda_\omega\hat\omega_s\right).
\]
Using \(Y=X\omega_*+\xi\), where \(\omega_*=\omega_*^{(t)}\) and
\(\xi=(\xi_1,\dots,\xi_{n_t})^\top\), we obtain
\[
\hat\omega_{\mathrm{TL}}-\omega_*
=
-\lambda_\omega A_n(\lambda_\omega)\delta
+
A_n(\lambda_\omega)\frac1{n_t}X^\top \xi,
\]
where
\[
A_n(\lambda):=(S_n+\lambda I_q)^{-1},
\qquad
\delta:=\omega_*-\hat\omega_s.
\]
This is the fundamental bias--variance decomposition.

\bigskip

\paragraph{Step 3: exact finite-sample conditional mean predictive risk.}

\begin{theorem}[Exact decomposition]
Under (A1)--(A6), for every fixed \(\lambda>0\),
\[
\bar r^{\mathrm{pred}}(\lambda\mid X)
=
\bar\tau_t^2
+
\lambda^2\,\mathrm{tr}\!\left\{
\Sigma_{RR}A_n(\lambda)\,\mathbb E[\delta\delta^\top]\,A_n(\lambda)
\right\}
+
\frac{\bar\tau_t^2}{n_t}\,
\mathrm{tr}\!\left\{
\Sigma_{RR}A_n(\lambda)S_nA_n(\lambda)
\right\},
\]
where
\[
\bar\tau_t^2:=\mathbb E[\tau_t^2(\beta^{(t)})]
=
\sigma_\varepsilon^2+\frac{\alpha_t^2}{p}\mathrm{tr}(\Sigma_{U\mid R}).
\]
\end{theorem}

\begin{proof}
For a fresh test point \((z_{\rm new},y_{\rm new})\), conditional on \(\beta^{(t)}\),
\[
y_{\rm new}=z_{\rm new}^\top\omega_*+\xi_{\rm new},
\]
with \(\xi_{\rm new}\) independent of \(z_{\rm new}\) and of the training sample, and
\(\mathbb E[\xi_{\rm new}^2\mid\beta^{(t)}]=\tau_t^2(\beta^{(t)})\).
Hence, conditional on \(X\),
\[
r^{\mathrm{pred}}(\hat\omega_{\mathrm{TL}})
=
\tau_t^2(\beta^{(t)})
+
(\hat\omega_{\mathrm{TL}}-\omega_*)^\top
\Sigma_{RR}
(\hat\omega_{\mathrm{TL}}-\omega_*).
\]
Taking the outer conditional expectation given \(X\) yields
\[
\bar r^{\mathrm{pred}}(\lambda\mid X)
=
\bar\tau_t^2
+
\mathbb E\!\left[
(\hat\omega_{\mathrm{TL}}-\omega_*)^\top
\Sigma_{RR}
(\hat\omega_{\mathrm{TL}}-\omega_*)
\,\middle|\,X
\right].
\]
Insert the identity
\[
\hat\omega_{\mathrm{TL}}-\omega_*
=
-\lambda A_n(\lambda)\delta
+
A_n(\lambda)\frac1{n_t}X^\top\xi.
\]
Because \(\mathbb E[\delta]=0\) and \(\delta\) is independent of \(\xi\) conditional on \(X\),
the cross term has mean zero. Therefore
\[
\mathbb E\!\left[
(\hat\omega_{\mathrm{TL}}-\omega_*)^\top
\Sigma_{RR}
(\hat\omega_{\mathrm{TL}}-\omega_*)
\,\middle|\,X
\right]
=
\lambda^2\mathrm{tr}\!\left\{
\Sigma_{RR}A_n(\lambda)\mathbb E[\delta\delta^\top]A_n(\lambda)
\right\}
+
\mathbb E\!\left[
\Bigl(\frac1{n_t}X^\top\xi\Bigr)^\top
A_n(\lambda)\Sigma_{RR}A_n(\lambda)
\Bigl(\frac1{n_t}X^\top\xi\Bigr)
\,\middle|\,X
\right].
\]
Now conditional on \(X\),
\[
\mathrm{Var}\!\left(\frac1{n_t}X^\top\xi\,\middle|\,X\right)
=
\frac1{n_t^2}X^\top\mathbb E[\xi\xi^\top]X.
\]
Since the \(\xi_i\) are independent, mean zero, and each has variance
\(\bar\tau_t^2\) after averaging over \(\beta^{(t)}\),
\[
\mathbb E[\xi\xi^\top]=\bar\tau_t^2 I_{n_t}.
\]
Hence
\[
\mathrm{Var}\!\left(\frac1{n_t}X^\top\xi\,\middle|\,X\right)
=
\frac{\bar\tau_t^2}{n_t}S_n.
\]
Using \(\mathbb E[v^\top Mv]=\mathrm{tr}(M\mathrm{Var}(v))\) for zero-mean \(v\),
\[
\mathbb E\!\left[
\Bigl(\frac1{n_t}X^\top\xi\Bigr)^\top
A_n(\lambda)\Sigma_{RR}A_n(\lambda)
\Bigl(\frac1{n_t}X^\top\xi\Bigr)
\,\middle|\,X
\right]
=
\frac{\bar\tau_t^2}{n_t}
\mathrm{tr}\!\left\{
\Sigma_{RR}A_n(\lambda)S_nA_n(\lambda)
\right\}.
\]
This proves the formula.
\end{proof}

\bigskip

\paragraph{Step 4: control of the transfer mismatch term \(\mathbb E[\delta\delta^\top]\).}

\begin{lemma}[Second-moment bound for the transfer mismatch]
Define
\[
\Delta_{ts}^2:=\alpha_t^2+\alpha_s^2-2\rho\alpha_t\alpha_s.
\]
Under (A3)--(A5),
\[
\mathbb E[\delta\delta^\top]
\preceq
\frac{\Delta_{ts}^2+C^U}{p}\,
\Phi^{(t)}(\Phi^{(t)})^\top
\preceq
\frac{(\Delta_{ts}^2+C^U)C_\Phi}{p}\,I_q.
\]
\end{lemma}

\begin{proof}
By definition,
\[
\delta
=
\omega_*-\hat\omega_s
=
\Phi^{(t)}\beta^{(t)}-\Phi^{(t)}\hat\beta^{(s)}
=
\Phi^{(t)}\Bigl[(\beta^{(t)}-\beta^{(s)})-e_s\Bigr].
\]
Since \(e_s\) is mean zero and independent of \((\beta^{(t)},\beta^{(s)})\),
\[
\mathbb E[\delta\delta^\top]
=
\Phi^{(t)}
\left\{
\mathrm{Var}(\beta^{(t)}-\beta^{(s)})+\mathbb E[e_s e_s^\top]
\right\}
(\Phi^{(t)})^\top.
\]
Now
\[
\mathrm{Var}(\beta^{(t)}-\beta^{(s)})
=
\frac{\alpha_t^2}{p}I_p+\frac{\alpha_s^2}{p}I_p
-
2\frac{\rho\alpha_t\alpha_s}{p}I_p
=
\frac{\Delta_{ts}^2}{p}I_p.
\]
By (A5),
\[
\mathbb E[e_s e_s^\top]\preceq \frac{C^U}{p}I_p.
\]
Therefore
\[
\mathbb E[\delta\delta^\top]
\preceq
\frac{\Delta_{ts}^2+C^U}{p}\,
\Phi^{(t)}(\Phi^{(t)})^\top.
\]
The second inequality follows from \(\|\Phi^{(t)}\|_{\mathrm{op}}^2\le C_\Phi\).
\end{proof}

\bigskip

\paragraph{Step 5: finite-sample upper bound.}

\begin{theorem}[Finite-sample conditional upper bound]
Under (A1)--(A6), for every \(\lambda>0\),
\[
\bar r^{\mathrm{pred}}(\lambda\mid X)
\le
\bar\tau_t^2
+
\frac{(\Delta_{ts}^2+C^U)C_\Phi}{p}\,
\lambda^2\,
\mathrm{tr}\!\left\{\Sigma_{RR}A_n(\lambda)^2\right\}
+
\frac{\bar\tau_t^2}{n_t}\,
\mathrm{tr}\!\left\{\Sigma_{RR}A_n(\lambda)S_nA_n(\lambda)\right\}.
\]
In particular, by \(\lambda_{\max}(\Sigma_{RR})\le C_\Sigma\),
\[
\bar r^{\mathrm{pred}}(\lambda\mid X)
\le
\bar\tau_t^2
+
C_\Sigma C_\Phi(\Delta_{ts}^2+C^U)\frac{\lambda^2}{p}\,
\mathrm{tr}\!\left\{A_n(\lambda)^2\right\}
+
C_\Sigma\bar\tau_t^2\frac1{n_t}\,
\mathrm{tr}\!\left\{A_n(\lambda)S_nA_n(\lambda)\right\}.
\]
\end{theorem}

\begin{proof}
Start from the exact decomposition in the previous theorem and use the matrix bound from the lemma:
\[
\lambda^2\mathrm{tr}\!\left\{
\Sigma_{RR}A_n(\lambda)\mathbb E[\delta\delta^\top]A_n(\lambda)
\right\}
\le
\frac{(\Delta_{ts}^2+C^U)C_\Phi}{p}\,
\lambda^2\,
\mathrm{tr}\!\left\{\Sigma_{RR}A_n(\lambda)^2\right\}.
\]
This gives the first inequality. The second follows from
\[
\Sigma_{RR}\preceq C_\Sigma I_q.
\]
\end{proof}

\bigskip

\paragraph{Step 6: asymptotic deterministic equivalent.}

Define for \(\lambda>0\),
\[
\frac1q\mathrm{tr}\{A_n(\lambda)^2\}\ \xrightarrow{\text{a.s.}}\ v_\eta(\lambda),
\qquad
\frac1q\mathrm{tr}\{A_n(\lambda)S_nA_n(\lambda)\}\ \xrightarrow{\text{a.s.}}\ u_\eta(\lambda).
\]
Using \(q/p\to\kappa\) and \(q/n_t\to\eta\), the finite-sample upper bound immediately yields:

\begin{theorem}[High-dimensional asymptotic upper bound]
Under (A1)--(A7), for every fixed \(\lambda>0\),
\[
\bar r^{\mathrm{pred}}(\lambda\mid X)
\le
\bar\tau_t^2
+
\kappa\,C_\Sigma C_\Phi(\Delta_{ts}^2+C^U)\,
\lambda^2 v_\eta(\lambda)
+
\eta\,C_\Sigma\bar\tau_t^2\,u_\eta(\lambda)
+
o_{\mathbb P}(1).
\]
Equivalently,
\[
\limsup_{n_t,p\to\infty}
\bar r^{\mathrm{pred}}(\lambda\mid X)
\le
\bar\tau_t^2
+
\kappa\,C_\Sigma C_\Phi(\Delta_{ts}^2+C^U)\,
\lambda^2 v_\eta(\lambda)
+
\eta\,C_\Sigma\bar\tau_t^2\,u_\eta(\lambda)
\]
in probability.
\end{theorem}

\begin{proof}
From the preceding theorem,
\[
\bar r^{\mathrm{pred}}(\lambda\mid X)
\le
\bar\tau_t^2
+
C_\Sigma C_\Phi(\Delta_{ts}^2+C^U)\frac{\lambda^2}{p}\,
\mathrm{tr}\{A_n(\lambda)^2\}
+
C_\Sigma\bar\tau_t^2\frac1{n_t}\,
\mathrm{tr}\{A_n(\lambda)S_nA_n(\lambda)\}.
\]
Now write
\[
\frac1p\mathrm{tr}\{A_n(\lambda)^2\}
=
\frac{q}{p}\cdot \frac1q\mathrm{tr}\{A_n(\lambda)^2\},
\qquad
\frac1{n_t}\mathrm{tr}\{A_n(\lambda)S_nA_n(\lambda)\}
=
\frac{q}{n_t}\cdot \frac1q\mathrm{tr}\{A_n(\lambda)S_nA_n(\lambda)\}.
\]
By (A1) and (A7),
\[
\frac1p\mathrm{tr}\{A_n(\lambda)^2\}
\to
\kappa v_\eta(\lambda),
\qquad
\frac1{n_t}\mathrm{tr}\{A_n(\lambda)S_nA_n(\lambda)\}
\to
\eta u_\eta(\lambda),
\]
in probability (indeed almost surely on the event of spectral convergence). Substituting these limits proves the claim.
\end{proof}

\bigskip

\paragraph{Oracle choice of \(\lambda_\omega\).}

The asymptotic upper bound is minimized by
\[
\lambda_\omega^{\mathrm{opt}}
\in
\arg\min_{\lambda>0}
\Bigl\{
\kappa\,C_\Sigma C_\Phi(\Delta_{ts}^2+C^U)\lambda^2 v_\eta(\lambda)
+
\eta\,C_\Sigma\bar\tau_t^2\,u_\eta(\lambda)
\Bigr\}.
\]
Since \(u_\eta(\lambda)\) is strictly decreasing in \(\lambda\) and
\(\lambda^2v_\eta(\lambda)\) is increasing for small \(\lambda\) and bounded as \(\lambda\to\infty\), the objective has the standard bias--variance tradeoff structure and admits at least one minimizer on \((0,\infty)\).

Differentiating,
\[
\frac{d}{d\lambda}\Bigl[\lambda^2 v_\eta(\lambda)\Bigr]
=
2\lambda v_\eta(\lambda)+\lambda^2 v_\eta'(\lambda),
\]
and
\[
u_\eta'(\lambda)
=
-2\int \frac{x}{(x+\lambda)^3}\,dF_\eta(x)<0.
\]
Hence any interior optimizer satisfies the first-order condition
\[
\kappa\,C_\Sigma C_\Phi(\Delta_{ts}^2+C^U)
\Bigl[2\lambda v_\eta(\lambda)+\lambda^2 v_\eta'(\lambda)\Bigr]
+
\eta\,C_\Sigma\bar\tau_t^2\,u_\eta'(\lambda)
=0.
\]
Equivalently,
\[
\kappa\,C_\Phi(\Delta_{ts}^2+C^U)
\Bigl[2\lambda v_\eta(\lambda)+\lambda^2 v_\eta'(\lambda)\Bigr]
+
\eta\,\bar\tau_t^2\,u_\eta'(\lambda)
=0.
\]
This equation characterizes the asymptotically optimal regularization level.

\bigskip

\paragraph{Transparent special case: \(\Sigma=I_p\).}

When \(\Sigma=I_p\), the situation becomes especially simple:
\[
\Sigma_{RR}=I_q,\qquad \Sigma_{RU}=0,\qquad \Phi^{(t)}=[I_q,0].
\]
Thus
\[
\omega_*^{(t)}=\beta_R^{(t)},
\qquad
\hat\omega_s^{(t)}=\hat\beta_R^{(s)},
\qquad
\Sigma_{U\mid R}=I_{p-q},
\]
and therefore
\[
\bar\tau_t^2
=
\sigma_\varepsilon^2+\alpha_t^2\frac{p-q}{p}
\ \longrightarrow\
\sigma_\varepsilon^2+\alpha_t^2(1-\kappa).
\]
Moreover \(C_\Sigma=C_\Phi=1\), and the asymptotic upper bound reduces to
\[
\bar r^{\mathrm{pred}}(\lambda\mid X)
\le
\bar\tau_t^2
+
\kappa(\Delta_{ts}^2+C^U)\lambda^2 v_\eta^{\mathrm{MP}}(\lambda)
+
\eta\bar\tau_t^2\,u_\eta^{\mathrm{MP}}(\lambda)
+
o_{\mathbb P}(1),
\]
where \(u_\eta^{\mathrm{MP}},v_\eta^{\mathrm{MP}}\) are computed under the Marchenko--Pastur law with aspect ratio \(\eta\). Thus the oracle tuning parameter is
\[
\lambda_{\omega,\mathrm{opt}}
\in
\arg\min_{\lambda>0}
\Bigl\{
\kappa(\Delta_{ts}^2+C^U)\lambda^2 v_\eta^{\mathrm{MP}}(\lambda)
+
\eta\bar\tau_t^2 u_\eta^{\mathrm{MP}}(\lambda)
\Bigr\}.
\]

\bigskip

\paragraph{Interpretation of the bound.}

The bound has three terms:
\[
\underbrace{\bar\tau_t^2}_{\text{irreducible error}}
\;+\;
\underbrace{
\kappa\,C_\Sigma C_\Phi(\Delta_{ts}^2+C^U)\lambda^2 v_\eta(\lambda)
}_{\text{transfer bias / mismatch}}
\;+\;
\underbrace{
\eta\,C_\Sigma\bar\tau_t^2 u_\eta(\lambda)
}_{\text{estimation variance}}.
\]
The quantity
\[
\Delta_{ts}^2=\alpha_t^2+\alpha_s^2-2\rho\alpha_t\alpha_s
\]
is the intrinsic target--source discrepancy. It is small when the source and target coefficients are strongly aligned (\(\rho\) large), and large when the two sites are poorly matched. The additional term \(C^U\) quantifies source-estimation uncertainty. Thus transfer is beneficial precisely when
\[
\Delta_{ts}^2+C^U
\]
is small enough that shrinking toward the source center decreases variance more than it introduces transfer bias.

\bigskip

\paragraph{Conclusion.}
Under the stated assumptions, the transfer-ridge estimator around the source-induced center \(\hat\omega_s^{(t)}=\Phi^{(t)}\hat\beta^{(s)}\) admits the exact conditional mean predictive-risk decomposition
\[
\bar r^{\mathrm{pred}}(\lambda\mid X)
=
\bar\tau_t^2
+
\lambda^2\,\mathrm{tr}\!\left\{
\Sigma_{RR}A_n(\lambda)\,\mathbb E[\delta\delta^\top]\,A_n(\lambda)
\right\}
+
\frac{\bar\tau_t^2}{n_t}\,
\mathrm{tr}\!\left\{
\Sigma_{RR}A_n(\lambda)S_nA_n(\lambda)
\right\},
\]
and hence the asymptotic upper bound
\[
\bar r^{\mathrm{pred}}(\lambda\mid X)
\le
\bar\tau_t^2
+
\kappa\,C_\Sigma C_\Phi(\Delta_{ts}^2+C^U)\lambda^2 v_\eta(\lambda)
+
\eta\,C_\Sigma\bar\tau_t^2u_\eta(\lambda)
+
o_{\mathbb P}(1).
\]
The asymptotically optimal regularization level is the minimizer of the deterministic equivalent on the right-hand side, equivalently the solution of the displayed first-order equation.

\end{document}
